\section{Completion and Synthesis Results}
\label{sec:phase-results}\label{sec:gap-study}
We first evaluate the distance phase independently of the mixer. This
isolates the effect of choosing unspecified table entries, then asks how
that effect depends on geometry, label structure and synthesis.

\subsection{Completion gains and synthesis dependence}
We find that completion improves on all four inexpensive extensions under
the same Gray synthesis method (Table~\ref{tab:extensions}). L1 uses fewer CX gates
than each alternative in all 96 cases with unused binary codes. All 64
power-of-two cases tie within every synthesis method, as required by the
uniqueness of their Walsh expansion. On the replication cohort, median
reductions relative to virtual-coordinate extension are 28.0\%, 53.9\% and
21.6\% at six, nine and twelve sites. The corresponding reductions relative
to nearest-code extension are 25.2\%, 53.0\% and 21.5\%.
Thus, on these geometries and label assignments, optimizing the extension
improves on natural alternatives to zero padding.

\begin{table}[t]\centering\small\setlength{\tabcolsep}{3pt}
\caption{Exact extensions under symbolic Gray synthesis. The initial
100-case and replication 60-case cohorts contain 20 and 12 geometries per
site count, respectively. CX counts are medians with interquartile ranges.
Percentage changes are medians of paired L1/reference changes, not ratios
of the displayed medians; negative values favor L1. Power-of-two controls
are omitted because all extensions coincide.}
\label{tab:extensions}\begin{tabular}{llrrrrr}
\toprule
Cohort & $m$ & L1 CX [IQR] & vs zero & vs virtual & vs nearest & vs mean\\
\midrule
100-case & 6 & 238 [236, 262] & -23.8\% & -26.3\% & -16.7\% & -23.8\%\\
100-case & 9 & 655 [619, 673] & -55.0\% & -55.0\% & -52.5\% & -55.0\%\\
100-case & 12 & 1095 [991, 1164] & -24.7\% & -23.0\% & -23.7\% & -24.7\%\\
60-case & 6 & 228 [204, 235] & -27.6\% & -28.0\% & -25.2\% & -27.6\%\\
60-case & 9 & 552 [438, 630] & -53.9\% & -53.9\% & -53.0\% & -53.9\%\\
60-case & 12 & 939 [858, 966] & -21.3\% & -21.6\% & -21.5\% & -21.3\%\\
\bottomrule
\end{tabular}

\end{table}

The benefit is sensitive to synthesis. With the angle fixed at
$\gamma=0.371$, generic diagonal synthesis reduces the replication-cohort
L1/virtual savings to 10.9\%, 1.3\% and 0.7\%. At these three sizes, L1
uses fewer/equal/more CX gates in 11/0/1, 8/4/0 and 11/1/0 cases,
respectively. L1 also improves on virtual extension throughout the
replication cases with unused codes under the tested PyZX procedure, but
that procedure produces more CX gates than numerical Gray synthesis for
every input (Appendix~\ref{sec:synthesis-diagnostics}). The large Gray
savings therefore describe a particular combination of completion and
synthesis, rather than a compiler-independent reduction.

We verify that these gate savings preserve the required action by
reconstructing all 1,760 stored coefficient vectors. The maximum
valid-distance residual is $5.69\times10^{-14}$ for the exact extensions. Replaying all 864 saved
confirmation and transfer parameter vectors with L1-completed token
energies gives the same energies and output probabilities at stored
floating-point precision. These checks support
Proposition~\ref{prop:completion}: the completion changes implementation
cost while preserving the ideal algorithm. Numerical angle conversion
and deliberate approximation have separate error checks in
Appendix~\ref{sec:synthesis-diagnostics}.


\subsection{The extent of completion freedom}
A two-register distance table has $2^{2k}$ entries but only $m^2$ prescribed
values. Its fraction of unspecified entries is therefore
\begin{equation}
 F=1-\frac{m^2}{2^{2\lceil\log_2m\rceil}}.
 \label{eq:free-fraction}
\end{equation}
This fraction concerns unused binary codes, not the number of vacant
placement sites. At a power of two, $F=0$ and completion cannot change the
operator. Between powers of two, freedom becomes available but does not
by itself determine the saving: six and twelve sites both have $F=7/16$,
yet their CX reductions differ (Fig.~\ref{fig:free-fraction}). Nine sites
have $F=175/256$. Geometry and label assignment determine the required
values, while net aggregation and parity synthesis determine how a chosen
extension translates into gates.

\begin{figure}[t]\centering
\includegraphics[width=.97\textwidth]{fig_reviewer_free_fraction.pdf}
\caption{CX reductions from L1 completion relative to zero extension.
The initial cohort uses fixed mask order (20 cases per size); the
replication cohort uses beam search and Gray synthesis (12 per size).
Boxes show medians and quartiles; whiskers reach the most extreme values
within 1.5 IQR, with other observations shown individually. Powers of two have no
completion freedom and hence zero reduction.}
\Description{Three panels show paired CX reductions by site count for fixed
mask order, beam search and Gray synthesis, with power-of-two controls at zero.}
\label{fig:free-fraction}
\end{figure}

\subsection{Joint optimization of coordinate contributions}
Manhattan distance is additive in its coordinates, suggesting a simpler
alternative to joint completion: solve Eq.~\eqref{eq:lp} separately for
$|x_a-x_b|$ and $|y_a-y_b|$, then add the coefficient vectors. The sum
$c_x+c_y$ is feasible for the joint problem. Consequently, the joint
optimum cannot have a larger weighted norm, although its synthesized
circuit need not have fewer gates. Both formulations retain the same
$O(m^4)$ dense constraint storage because both use the full site-label
domain. Coordinate separation alone does not solve the storage problem.

\begin{table}[t]\centering\small
\caption{Coordinate-separated relative to joint L1 completion. Each row
contains 20 initial-cohort or 12 replication-cohort cases. W/T/L counts
cases with fewer/equal/more CX gates for the separated construction.
Changes and IQRs summarize paired percentages. Norm comparisons in the
text evaluate the LP objective on $c_x+c_y$. Power-of-two ties are omitted.}
\label{tab:coordinate}\begin{tabular}{llrrrl}
\toprule
Cohort & Synthesis & $m$ & CX change & IQR & W/T/L\\
\midrule
100-case & fixed order & 6 & +6.6\% & [+0.0, +26.9] & 0/9/11\\
100-case & fixed order & 9 & +12.1\% & [+0.0, +27.7] & 1/6/13\\
100-case & fixed order & 12 & +21.6\% & [+0.0, +36.1] & 0/6/14\\
60-case & beam search & 6 & +11.6\% & [+0.0, +22.1] & 1/3/8\\
60-case & beam search & 9 & +17.0\% & [+6.6, +27.7] & 0/3/9\\
60-case & beam search & 12 & +17.1\% & [+7.2, +25.7] & 0/3/9\\
60-case & Gray & 6 & +10.0\% & [+0.0, +19.0] & 1/3/8\\
60-case & Gray & 9 & +16.3\% & [+2.1, +28.9] & 0/3/9\\
60-case & Gray & 12 & +14.6\% & [+4.9, +18.8] & 0/3/9\\
\bottomrule
\end{tabular}

\end{table}

Joint completion has a strictly smaller weighted norm in 67 of the 96
cases with unused codes. Coordinate-separated completion increases
replication-cohort Gray CX by median 10.0\%, 16.3\% and 14.6\% at six,
nine and twelve sites (Table~\ref{tab:coordinate}). Some individual cases
nevertheless favor separation in gate count, illustrating the distinction
between the LP objective and the synthesized circuit. All 64 power-of-two
cases yield identical coefficients, and the maximum valid-distance
residual is $8.35\times10^{-14}$. Joint optimization can exploit
cancellation between coordinate contributions that separate solves miss.

\subsection{Direct construction on regular arrays}
The recurrence in Section~\ref{sec:structured} takes a different approach:
it uses coordinate bits directly and avoids the dense LP. We evaluate it
on binary rows and row-major rectangles, including prefixes with unused
codes, at the same five site counts. Each comparison uses a four-cell,
six-net phase and Gray synthesis (Table~\ref{tab:structured}). These deterministic geometries examine
the construction's mechanism; they are separate from the shuffled-label
cohorts.

\begin{table}[t]\centering\small
\caption{Gray CX counts for regular-coordinate distance phases with six
nets. Each row is one deterministic geometry. All three extensions are
exact on valid pairs; the recurrence assigns distances to unused binary
coordinates. These are construction examples rather than a sampled cohort.}
\label{tab:structured}\begin{tabular}{lrrrr}
\toprule
Coordinate set & $m$ & Zero CX & L1 CX & Structured CX\\
\midrule
row & 6 & 292 & 238 & 264\\
row & 8 & 264 & 264 & 264\\
row & 9 & 1406 & 344 & 720\\
row & 12 & 1286 & 706 & 720\\
row & 16 & 720 & 720 & 720\\
grid & 6 & 318 & 86 & 86\\
grid & 8 & 86 & 86 & 86\\
grid & 9 & 1262 & 110 & 148\\
grid & 12 & 1182 & 110 & 148\\
grid & 16 & 148 & 148 & 148\\
\bottomrule
\end{tabular}

\end{table}

The recurrence reproduces every coordinate-pair distance through six bits,
and its support count agrees with Eq.~\eqref{eq:structured-support} through
sixteen bits. The latter check constructs coefficients rather than a
placement circuit at that width. At six grid sites, the recurrence and L1
both use 86 CX, compared with 318 for zero extension. The recurrence is not
always as economical as the LP: at nine grid sites it uses 148 CX versus
110, and at nine row sites 720 versus 344. It therefore provides a more
compact construction on this structured subclass without necessarily
minimizing the resulting circuit.
